\subsection{The survey}
\label{data:survey}

The data come from a cross-sectional household survey: each respondent was
interviewed exactly once, in person, in Mandarin Chinese. Interviews were
conducted in two counties bordering Giant Panda National Park, Sichuan
Province -- Baoxing (15--31 August 2024) and Qingchuan (10 March--10 April
2025) -- seven months apart. The analysis data set supplied by the project partner contains 625
respondents and 57 columns (an identifier, three location labels, and 53
survey variables): 160 respondents from Baoxing and 465 from Qingchuan. Because
the two counties were not surveyed concurrently, place of residence and survey
timing cannot be statistically separated from each other in anything that
follows; every model in this report treats them as a single combined
\emph{site} covariate for that reason.

The questionnaire covered five thematic blocks: (1) sociodemographics (age,
gender, education, annual income, household assets, cropland area); (2)
conflict experience (frequency, duration, severity, species diversity, and
years of exposure to wild boar and other wildlife); (3) economic and
intangible costs (estimated damaged crop area and economic loss, alongside a
general attitude rating and a free-text three-word description of wild boar);
(4) compensation and mitigation (application for compensation, satisfaction
with the amount received, mitigation strategies used, and willingness to pay
for wildlife-damage insurance); and (5) tolerance, measured with four Likert
items described in Section~\ref{measurement:tolerance}.

\subsection{Compensation status is confounded with county}
\label{data:confound}

The central grouping variable throughout this project is \emph{compensation
participation} (construct code \emph{CE0}): a binary indicator of whether a
household applied for compensation for wild boar damage (277 applicants, 348
non-applicants). Before any model could treat CE0 as a meaningful predictor, we
checked whether it was confounded with the county/survey-wave split described
above. It is, strongly: 119 of 160 respondents in Baoxing (74\%) reported
applying for compensation, versus only 158 of 465 in Qingchuan (34\%). A
Pearson chi-squared test of independence on the resulting $2\times2$ table
(county $\times$ applied for compensation) confirms this gap is far larger than
sampling noise would produce. For each cell $i$ of the table, let $O_i$ be
the observed count and $E_i$ the count expected under independence (row
total times column total, divided by the grand total); the test statistic

\begin{equation}
\chi^2 = \sum_{i} \frac{(O_i - E_i)^2}{E_i}
\label{eq:chisq}
\end{equation}

is compared against a $\chi^2$ distribution with $df = (r-1)(c-1) = 1$
degree of freedom for an $r=c=2$ table ($\chi^2 = 77.1$, $p < .001$, Yates'
continuity-corrected; uncorrected $\chi^2 = 78.7$; Fisher's exact test
$p = 5.8\times10^{-19}$). The corresponding odds ratio is 5.64: a Baoxing
resident had 5.6 times the odds of having applied for compensation, compared
with a Qingchuan resident. Because county and survey wave move together with
compensation status, every model in Sections~\ref{model} and
\ref{robustness} that includes CE0 as a predictor also includes site as an
explicit control.

\subsection{Compensation recipients differ in assets, not in who they are}
\label{data:balance}

We also compared the two compensation groups on the sociodemographic and
asset variables available for the full sample, using a Wilcoxon rank-sum
(Mann-Whitney U) test for each continuous variable and a chi-squared test for
the one categorical variable, gender. This is a covariate-balance check in
the spirit of a baseline comparison table in an observational study: it asks
whether recipients and non-recipients of compensation are otherwise similar
households, not whether these variables predict compensation uptake in a
multivariate sense. Pooled over both counties, recipients have significantly more total
cropland, more total livestock, and more household assets (all $p < .001$),
and a higher share of income from farming ($p = .017$), whereas age
($p = .401$), gender ($p = .318$), education ($p = .383$), and annual income
($p = .599$) show no significant difference between the two groups. Read
on its own this would suggest that compensation is claimed against exposure.
The next two paragraphs show that this reading is only partly supported,
because these pooled tests do not account for county.

A direction-of-analysis clarification is worth making explicit here, since
it is easy to read the comparison above the wrong way round. The tests just
reported ask, for a variable already split into two known groups (applied
for compensation vs.\ not), whether that variable's \emph{distribution}
differs between the groups -- for example, "do cropland holdings differ
between applicants and non-applicants?" They do not ask the reverse,
predictive-sounding question -- "if a respondent is male, is he more likely
to apply for compensation?" -- which would instead treat compensation status
as an outcome to be predicted \emph{from} gender, age, and the other
variables. Both questions are legitimate, but they are different questions,
and this balance check deliberately asks the first one: its purpose is to
establish whether the two compensation groups are otherwise comparable
before any later comparison between them is trusted, in the same way a
baseline-characteristics table works in any observational study. It is not
an attempt at a predictive model of who applies for compensation, and should
not be read as one.

The second, predictive-direction question does have an answer in this
report -- it is simply asked separately, inside the main structural model,
rather than at this bivariate, exploratory stage. The full model
(Section~\ref{model:spec}) includes explicit paths from every
sociodemographic variable \emph{to} compensation participation (CE0), estimated
jointly with every other path in the model. Table~\ref{tab:demog-to-ce0}
reports them: household income is the one variable that significantly
predicts compensation participation ($\beta = 0.080$, $p = .047$); gender is
close to but does not cross the conventional threshold ($\beta = 0.069$,
$p = .060$); age, education, and household assets do not reach significance.
So, to answer the specific example directly: being male is associated with a
somewhat higher chance of having applied for compensation, but that
association is borderline, not statistically significant at the conventional
5\% level, once the other background variables are accounted for
simultaneously.

\paragraph{Checking the pooled Wilcoxon results in the predictive direction.}
Because the rank-sum tests above are bivariate and pooled over two very
different counties, we checked them in two further ways
(Table~\ref{tab:balance-check}). First, we repeated each test within each
county. The pooled differences in livestock and household assets disappear
in Qingchuan ($p = .51$ and $p = .75$) and are present only in Baoxing, so
they largely reflect county, not a general link between farm size and
applying. Second, we fitted a logistic regression of compensation
participation on these variables and gender together
(Table~\ref{tab:balance-check}). Writing $\pi_i$ for the probability that
household $i$ applied, the model without county is
\begin{equation}
\begin{split}
\log\frac{\pi_i}{1-\pi_i} = \beta_0 &+ \beta_1\,\text{age}_i + \beta_2\,\text{male}_i
 + \beta_3\,\text{education}_i + \beta_4\,\text{household size}_i
 + \beta_5\,\text{income}_i \\ &+ \beta_6\,\text{agricultural share}_i
 + \beta_7\,\text{cropland}_i + \beta_8\,\text{livestock}_i
 + \beta_9\,\text{assets}_i ,
\end{split}
\label{eq:logit}
\end{equation}
and the model with county adds the term $\beta_{10}\,\text{Baoxing}_i$,
an indicator that is 1 for Baoxing and 0 for Qingchuan. Income, cropland,
livestock, and assets enter as $\log(1+x)$; all predictors except the two
indicators are standardised, so $\exp(\beta_j)$ is the odds ratio for a
one-standard-deviation increase.
Without county, livestock is the clearest predictor (odds ratio 1.63 per
standard deviation, $p < .001$). Once county is added, livestock is no longer
significant, county itself is by far the strongest predictor (odds ratio 8.7),
and cropland (odds ratio 1.66, $p < .001$) is clearly positively associated
with applying. Income (odds ratio 1.26, $p = .017$) and household assets
(odds ratio 0.65, $p = .010$) point in opposite directions and are only
borderline once the number of tests is accounted for (Table~\ref{tab:test-count});
the opposite signs most likely reflect collinearity among the three wealth
measures, so we do not interpret them one by one. The consistent
conclusion is narrower than our first reading: county is the dominant
difference between applicants and non-applicants, and, within county, larger
cropland holdings go with a higher chance of having applied.

\paragraph{How many tests this rests on, and what we finally claim.}
The comparison of the two compensation groups involved several rounds of
testing, and we state them all so that the reader can judge how much weight a
single $p$-value deserves (Table~\ref{tab:test-count}). The first round, the
balance table, consists of 10 tests: eight Wilcoxon rank-sum tests, one
chi-squared test for gender, and one for county. The within-county repetition
adds 16 Wilcoxon tests (eight variables in each of two counties), and the two
logistic regressions add 19 coefficient tests (nine without county, ten with
it), 45 tests in all. The five paths from background variables to compensation
participation in the structural model are part of the main model and are not
counted here. None of the tests was adjusted for multiplicity in the original
balance table. With 10 tests at the 5\% level, the chance of at least one
false positive would be up to about 40\% if no variable truly differed, so
we applied two standard corrections within each round. The Holm procedure
\citep{Holm1979} controls the probability of making \emph{even one} false
positive among the $m$ tests of a round: the $p$-values are sorted from
smallest to largest, the smallest is multiplied by $m$, the second by $m-1$,
and so on, and a result counts as significant only if its adjusted value and
all smaller ones are below .05. The Benjamini--Hochberg procedure
\citep{Benjamini1995} controls the less strict \emph{false discovery rate},
the expected share of false positives among the results declared
significant: the $i$-th smallest $p$-value is multiplied by $m/i$, giving an
adjusted value $q$. Holm suits a setting where a single false claim is
costly; Benjamini--Hochberg suits a screening setting such as this one,
where we look at many variables and accept that a small share of the flagged
ones may be chance. We use Benjamini--Hochberg as the main criterion and
report Holm where the two disagree. For example, the share of income from
agriculture has the fifth-smallest of the 10 $p$-values in the first round
($p = .017$): its Benjamini--Hochberg value is $.017 \times 10/5 = .034$, and
its Holm value is $.017 \times 6 = .10$. This does not change the
main conclusions: in the first round, cropland, livestock, assets, and county
remain significant, and the share of income from agriculture, which was only
nominally significant ($p = .017$), remains so at a 5\% false-discovery rate
($q = .034$) but not under the stricter Holm correction ($p = .10$); we
therefore give it little weight. In the logistic regression with county,
cropland ($q < .001$) and county ($q < .001$) are firm, while assets
($q = .033$) and income ($q = .041$) are borderline.

Our final claims about who applies for compensation are therefore these.
\begin{enumerate}
  \item County is by far the largest difference between applicants and
    non-applicants (odds ratio 8.7 in the joint model; $\chi^2 = 77.1$).
  \item Larger cropland holdings go with a higher chance of having applied,
    both within each county and in the joint model. This is the only
    farm-related variable for which the evidence is consistent.
  \item The pooled differences in livestock and household assets are mostly
    a county effect: they are absent in Qingchuan and present only in Baoxing.
    We therefore do \emph{not} claim that compensation tracks livestock or
    asset exposure in general.
  \item Age, gender, and education show no reliable difference between the
    groups, and annual income and the share of income from agriculture none
    that we would interpret.
  \item Compensation status is therefore not standing in for wealth or
    demographics in general, but we cannot separate it from county and, to a
    lesser extent, from the size of a household's cropland, which is why both
    are kept as controls in the main model.
\end{enumerate}

\begin{table}[htbp]
\centering
\caption{Tests behind the comparison of compensation groups. Significant
means significant at 5\% after Benjamini--Hochberg correction within the
round.}
\label{tab:test-count}
\small
\begin{tabular}{p{4.8cm}rrp{4.4cm}}
\hline
\textbf{Round} & \textbf{Tests} & \textbf{Significant} & \textbf{Which} \\
\hline
Balance table (pooled): 8 Wilcoxon, 2 chi-squared & 10 & 5 & Cropland, livestock, assets, share of agricultural income, county \\
Wilcoxon within Qingchuan & 8 & 2 & Education, cropland (both $q = .043$) \\
Wilcoxon within Baoxing & 8 & 5 & Age, income, cropland, livestock, assets \\
Logistic regression without county & 9 & 1 & Livestock \\
Logistic regression with county & 10 & 4 & County, cropland; assets and income borderline \\
\hline
\textbf{Total} & \textbf{45} & & \\
\hline
\end{tabular}
\end{table}

\begin{table}[htbp]
\centering
\caption{Who applies for compensation? Pooled and within-county Wilcoxon
rank-sum $p$-values, and odds ratios (per standard deviation; money variables
log-transformed) from a logistic regression of compensation participation
on all variables jointly, including county, with heteroscedasticity-robust
confidence intervals.}
\label{tab:balance-check}
\small
\begin{tabular}{p{4.6cm}rrrr}
\hline
 & \multicolumn{3}{c}{\textbf{Wilcoxon $p$}} & \textbf{Odds ratio} \\
\textbf{Variable} & Pooled & Qingchuan & Baoxing & [95\% CI] \\
\hline
Age                      & .401   & .884 & .002 & 0.94 [0.76, 1.15] \\
Education                & .383   & .019 & .091 & 1.22 [1.00, 1.50] \\
Household size           & .270   & .157 & .552 & 0.82 [0.68, 0.99] \\
Annual income            & .599   & .567 & .006 & 1.26 [1.04, 1.53] \\
Share of income from agriculture & .017 & .274 & .306 & 1.04 [0.85, 1.27] \\
Total cropland           & $<.001$ & .017 & $<.001$ & 1.66 [1.28, 2.16] \\
Total livestock          & $<.001$ & .508 & $<.001$ & 1.21 [0.96, 1.53] \\
Household assets         & $<.001$ & .746 & .004 & 0.65 [0.47, 0.90] \\
Gender (male)            & --- & --- & --- & 1.40 [0.93, 2.10] \\
County (Baoxing)         & --- & --- & --- & 8.68 [4.94, 15.2] \\
\hline
\end{tabular}
\end{table}

\begin{table}[htbp]
\centering
\caption{Sociodemographic paths to compensation participation (CE0), from the
full structural model (Section~\ref{model:spec}), estimated jointly with
every other path in the model.}
\label{tab:demog-to-ce0}
\begin{tabular}{lrr}
\hline
\textbf{Predictor} & \textbf{$\beta$} & \textbf{$p$} \\
\hline
Gender     & 0.069  & .060 \\
Age        & $-0.017$ & .690 \\
Education  & 0.069  & .100 \\
Income (log)     & 0.080  & .047 \\
Household assets (log) & $-0.006$ & .900 \\
\hline
\end{tabular}
\end{table}

\subsection{Skewness, missing values, and outliers}
\label{data:quality}

Several monetary and area variables -- insurance willingness-to-pay, income,
household assets, and damage-related economic loss -- are heavily right-skewed
(household assets alone range from 77 to 3{,}087{,}729 RMB). All four are
log(1+x)-transformed before entering any model, so that a handful of very
large landholders do not dominate otherwise linear relationships.

Is this transformation necessary? For our conclusions, it is not. PLS-SEM
makes no normality assumption (Section~\ref{model:why-plssem}), so skewness
is not a modelling requirement here. We therefore re-ran the observed-score
cross-check regression of tolerance (Section~\ref{app:transform}) with the
money variables left raw, rank-transformed, and log-transformed.
The coefficients that carry our conclusions are essentially identical in all
four versions: compensation participation $-0.165$ to $-0.168$, intangible
cost $-0.29$ to $-0.30$, intangible benefit $+0.11$ (Table~\ref{tab:transform}).
Only a borderline coefficient moves: household assets, whose $p$-value is
$.026$ with the log transform and $.064$ raw, a non-significant coefficient
in the structural model in either case.

The transformation was a precaution against a few extreme households
(assets up to 3 million RMB, against a typical range of a few thousand)
influencing coefficients, decided from each variable's raw distribution
before any model was fit and applied uniformly to all four variables. It is
a harmless choice, not an important one, and the results would be the same
without it. It is also imperfect: because income has zeros, $\log(1+x)$
turns its right skew (skewness $3.0$) into a left skew ($-3.7$).

Missingness in this dataset is not uniform, and treating all of it the same
way would have been a mistake. Two compensation-experience items (timeliness
and payment satisfaction, columns CE1 and CE2) are missing for the 348
respondents who never applied for compensation -- by survey design (skip
logic), not data loss -- and are therefore excluded from the main structural
model entirely, used only in supplementary applicant-only checks. Nine
respondents' lowest-accepted-percentage answers and four respondents'
expected-percentage answers (columns CA1, CA2) exceed 100\% of assessed loss
(up to 800\%); these are kept as genuine responses rather than capped or
deleted, since a respondent can coherently report wanting more than their
assessed loss, and a winsorised (capped-at-100\%) version was run separately
as a sensitivity check with essentially unchanged results. Zero-valued
responses that are legitimate answers -- prevention effectiveness, channel
count, and compensation-anticipation variables that define zero as "used
nothing" or "expects nothing" -- were likewise kept untouched. Apart from
the two skip-logic columns, and one respondent whose village label is
missing (the village label is not used in any model), the dataset has no
missing values in the other variables for all 625 respondents.

Finally, we screened for multivariate outliers using Mahalanobis distance
across the 15 core indicators of the Wildlife Tolerance Model (seven conflict-
experience items, two cost items, two benefit items, and four tolerance
items). For respondent $i$ with indicator vector $x_i \in \mathbb{R}^{15}$,
sample mean vector $\bar{x}$, and sample covariance matrix $S$, the squared
Mahalanobis distance is

\begin{equation}
D_i^2 = (x_i - \bar{x})^\top S^{-1} (x_i - \bar{x}).
\label{eq:mahalanobis}
\end{equation}

The matrix $S^{-1}$ rescales and decorrelates the 15 indicators before the
squared distance is taken, which is what lets $D_i^2$ account for how the
indicators correlate with each other rather than treating each one's
distance from the mean independently. Under multivariate normality, $D_i^2$
approximately follows a chi-squared distribution with 15 degrees of freedom
(one per variable, by the same logic as Equation~\ref{eq:chisq}); we
flagged respondents whose distance exceeded the $p < .001$ cutoff
($D^2 > 37.7$). Fourteen of 625 respondents (2.2\%) were
flagged -- split evenly across county and compensation status (seven from
each county, seven with and seven without compensation), so the flagged
cases are not concentrated in any one group. These respondents were kept in
the primary analysis (a statistical outlier is not automatically a data
error), and every core structural result reported in Section~\ref{model} was
re-estimated excluding them as a robustness check: all paths were stable in
sign, magnitude, and significance.
